\documentclass[a4paper, 11pt]{book}

% ─── Encoding & Language ─────────────────────────────────────────
\usepackage[utf8]{inputenc}
\usepackage[english]{babel}
\usepackage{microtype}

% ─── Layout ──────────────────────────────────────────────────────
\usepackage{geometry}
\geometry{left=2.5cm, right=2.5cm, top=2.8cm, bottom=2.8cm}

% ─── Fonts ───────────────────────────────────────────────────────
\usepackage{lmodern}
\usepackage[T1]{fontenc}
\usepackage{inconsolata}
\renewcommand{\familydefault}{\ttdefault}

% ─── Colour ──────────────────────────────────────────────────────
\usepackage[dvipsnames,svgnames,x11names]{xcolor}

\definecolor{forthDark}{HTML}{0B1B0B}
\definecolor{forthBlue}{HTML}{2ECC40}
\definecolor{forthCyan}{HTML}{7DCEA0}
\definecolor{forthGold}{HTML}{CCAA00}
\definecolor{forthOrange}{HTML}{CC7700}
\definecolor{forthBg}{HTML}{0D1B0D}
\definecolor{forthGrey}{HTML}{6B8E6B}
\definecolor{forthGreen}{HTML}{00FF41}

% ─── Graphics & Tables ──────────────────────────────────────────
\usepackage{graphicx}
\usepackage{wrapfig}
\usepackage{multirow}
\usepackage{booktabs}

% ─── Lists ───────────────────────────────────────────────────────
\usepackage{enumitem}
\setlist[enumerate]{label=\textcolor{forthBlue}{\roman*.}, leftmargin=*, itemsep=4pt}
\setlist[itemize]{label=\textcolor{forthBlue}{\textbullet}, leftmargin=*, itemsep=4pt}

% ─── Boxes ───────────────────────────────────────────────────────
\usepackage[most]{tcolorbox}

\usepackage{listings}
\lstset{
  basicstyle=\ttfamily\small,
  breaklines=true,
  columns=fixed,
  keepspaces=true,
  aboveskip=0pt,
  belowskip=0pt,
}

\newtcblisting{forthbox}[1][]{%
  enhanced,
  breakable,
  colback=forthBg,
  colframe=forthBlue,
  coltext=forthGreen,
  coltitle=white,
  fonttitle=\bfseries\sffamily,
  boxrule=0.6pt,
  arc=3pt,
  left=6mm, right=6mm, top=4mm, bottom=4mm,
  attach boxed title to top left={yshift=-3mm, xshift=8mm},
  boxed title style={
    colback=forthBlue,
    arc=2pt,
    boxrule=0pt,
    left=4mm, right=4mm, top=1.5mm, bottom=1.5mm
  },
  title={#1},
  shadow={1mm}{-1mm}{0mm}{black!8},
  listing only,
  listing options={basicstyle=\ttfamily\small\color{forthGreen}, breaklines=true, columns=fixed, keepspaces=true},
}

\newtcolorbox{notebox}{%
  enhanced,
  breakable,
  colback=forthGold!8,
  colframe=forthGold,
  boxrule=0pt,
  borderline west={3pt}{0pt}{forthGold},
  arc=0pt,
  left=5mm, right=5mm, top=3mm, bottom=3mm,
  fontupper=\sffamily\small,
}

% ─── Section styling ────────────────────────────────────────────
\usepackage{titlesec}

\titleformat{\chapter}[display]
  {\normalfont\huge\bfseries\sffamily\color{forthDark}}
  {\tikz[remember picture,overlay]{%
    \fill[forthBlue!10] (-1,-1) rectangle (\textwidth+1,2.2);
    \node[anchor=west, text=forthBlue!30, font=\fontsize{72}{72}\selectfont\bfseries\sffamily]
      at (-0.5,0.4) {\thechapter};
  }}
  {0pt}
  {\Huge}
  [\vspace{-4pt}{\color{forthBlue}\rule{\textwidth}{1.5pt}}]

\titleformat{\section}
  {\normalfont\Large\bfseries\sffamily\color{forthDark}}
  {\textcolor{forthBlue}{\thesection}}{0.6em}{}
  [\vspace{-6pt}{\color{forthBlue!30}\rule{\textwidth}{0.5pt}}]

\titleformat{\subsection}
  {\normalfont\large\bfseries\sffamily\color{forthDark}}
  {\textcolor{forthCyan}{\thesubsection}}{0.6em}{}

\titlespacing*{\chapter}{0pt}{-20pt}{30pt}
\titlespacing*{\section}{0pt}{18pt}{10pt}
\titlespacing*{\subsection}{0pt}{14pt}{8pt}

% ─── Headers & Footers ──────────────────────────────────────────
\usepackage{fancyhdr}
\pagestyle{fancy}
\fancyhf{}
\fancyhead[LE]{\sffamily\small\textcolor{forthGrey}{\leftmark}}
\fancyhead[RO]{\sffamily\small\textcolor{forthGrey}{\rightmark}}
\fancyfoot[C]{\sffamily\small\textcolor{forthBlue}{\thepage}}
\renewcommand{\headrulewidth}{0.4pt}
\renewcommand{\headrule}{\hbox to\headwidth{\color{forthBlue!30}\leaders\hrule height \headrulewidth\hfill}}
\renewcommand{\footrulewidth}{0pt}

\fancypagestyle{plain}{%
  \fancyhf{}
  \fancyfoot[C]{\sffamily\small\textcolor{forthBlue}{\thepage}}
  \renewcommand{\headrulewidth}{0pt}
}

% ─── Hyperlinks ──────────────────────────────────────────────────
\usepackage{hyperref}
\hypersetup{
  colorlinks=true,
  linkcolor=blue,
  urlcolor=forthOrange,
  citecolor=forthGreen,
  pdfauthor={Henry Hudson},
  pdftitle={Henceforth},
}

% ─── TikZ ────────────────────────────────────────────────────────
\usepackage{tikz}
\usetikzlibrary{calc}

% ─── Misc ────────────────────────────────────────────────────────
\usepackage{amsmath}
\usepackage{amssymb}
\usepackage{pifont}
\newcommand{\cmark}{\textcolor{forthBlue}{\ding{51}}}%
\newcommand{\xmark}{\textcolor{forthGrey}{\ding{55}}}%
\usepackage{fancyvrb}

\setlength{\parindent}{0pt}
\setlength{\parskip}{6pt plus 2pt minus 1pt}

\let\cleardoublepage\clearpage

% ═════════════════════════════════════════════════════════════════
\begin{document}

% ─── Title Page ──────────────────────────────────────────────────
\begin{titlepage}
\begin{tikzpicture}[remember picture, overlay]
  % ── Background ──
  \fill[black] (current page.south west) rectangle (current page.north east);
  % ── Accent line across the page ──
  \fill[forthOrange] ($(current page.west)+(0,0.5)$) rectangle ($(current page.east)+(0,0.58)$);
  % ── Title ──
  \node[anchor=west, text=white, font=\fontsize{52}{56}\selectfont\bfseries\ttfamily]
    at ($(current page.west)+(2.5,3.5)$) {Henceforth};
  % ── Subtitle ──
  \node[anchor=west, text=forthCyan, font=\Large\ttfamily]
    at ($(current page.west)+(2.5,1.8)$) {A FORTH Programming Environment for iOS};
  % ── Author ──
  \node[anchor=west, text=forthGrey, font=\large\ttfamily]
    at ($(current page.west)+(2.5,-0.8)$) {Henry Hudson};
  % ── Version / year ──
  \node[anchor=south east, text=forthGrey!60, font=\small\ttfamily]
    at ($(current page.south east)+(-2.5,1.5)$) {Documentation \(\cdot\) 2026};
  % ── Decorative prompt at bottom left ──
  \node[anchor=south west, text=forthBlue, font=\ttfamily\Large, opacity=0.4]
    at ($(current page.south west)+(2.5,1.5)$) {ok.};
\end{tikzpicture}
\end{titlepage}

% ─── Front Matter ────────────────────────────────────────────────
\pagenumbering{roman}
\tableofcontents
\newpage
\pagenumbering{arabic}

% ═════════════════════════════════════════════════════════════════
\chapter{About}
% ═════════════════════════════════════════════════════════════════

\section{Links to Source Code}

\href{https://github.com/henryhudson/FORTHapp}{github.com/henryhudson/FORTHapp}

Running the project requires Xcode on any iPhone / iPad device. Will soon be optimised for all Apple products.

Private for the time being --- get in touch if you would like access.

\section{Support}

\href{https://x.com/forth_hence4h}{Follow us on Twitter}

% ═════════════════════════════════════════════════════════════════
\chapter{Goals}
% ═════════════════════════════════════════════════════════════════

\section{Integrate with Bitcoin (BSV)}

\begin{enumerate}
\item Fetch latest block headers
\item Integrate with SwiftBSV for transactions, verification, wallet functions etc.
\item OPCodes file available to load and make use of
\item Make use of Timestamp server for users words and CRDT (Conflict-free Replicated Data Type) system
\item Full sCrypt capability --- each FORTH file will create a sCrypt that runs the op codes within
\end{enumerate}

\section{Get People Coding}

Provide a simple interface for the FORTH programming language, making it easy for people to learn to code.

\section{Access to Files App}

If permission is granted by the user, access to their documents will allow the program to load \texttt{.fs}, \texttt{.forth} and possibly other files for the user to use how they wish.

Create, read, update, delete --- CRUD operations should be supported, at least for FORTH files.

If access is granted, files will be stored on user's iCloud for backup/safety.

\section{Further Maths}

Eventually add functionality for trigonometry and calculus.

\section{Settings}

Colour customisation is available for users to pick a foreground colour and a background colour. Also allow users to pick between the various keyboards.

% ═════════════════════════════════════════════════════════════════
\chapter{Credits}
% ═════════════════════════════════════════════════════════════════

\section{Starting Forth --- Leo Brodie}

Leo Brodie, author of \emph{Starting Forth} --- a great textbook, suitable for both novices and professionals.

\href{https://1scyem2bunjw1ghzsf1cjwwn-wpengine.netdna-ssl.com/wp-content/uploads/2018/01/Starting-FORTH.pdf}{Starting FORTH (PDF)}

\section{BITSRFR}

Fellow surfer of the binary ocean who has made FORTH tutorials and is making a Swift command line FORTH application.

\href{https://www.youtube.com/user/JoeKreydt}{YouTube} \(\cdot\) \href{https://github.com/josephkreydt/cruxForth}{cruxFORTH on GitHub}

\section{Bitcoin}

The greatest monetary system ever created. Invented by Craig Wright aka Satoshi Nakamoto --- please read the white paper.

\href{https://craigwright.net/bitcoin-white-paper.pdf}{Bitcoin White Paper}

% ═════════════════════════════════════════════════════════════════
\chapter{Documentation}
% ═════════════════════════════════════════════════════════════════

Everything in FORTH is a \textbf{word} separated by spaces. That's all there is to the syntax. It's so simple.

\begin{notebox}
\textbf{FORTH is LIFO} --- FORTH is a stack-based programming language which operates on a \emph{Last In, First Out} principle.
\end{notebox}

\section{Built-in Words}

Ideally we give the user just enough to make use of FORTH and through FORTH some of Swift's functionality such as networking requests, the timer built into Swift, etc. But not so much as to be relied upon --- user creation of words should be encouraged, and if possible a community database of words and their definitions be made.

% ──────────────────────────────────────────────────────────────────
\subsection{Compile Words}

One of the great things about FORTH is its ability to extend the compiler during run time.

\begin{forthbox}[Compile Operators]
  :           puts the user into compile mode where they can
              start to define their own word.
  ;           exits compile mode so that word interpretation
              can again occur.
  recurse     calls the word being executed.
  exit        exits the current word.
  immediate   marks the most recently defined word as immediate,
              causing it to execute during compilation rather
              than being compiled.

  postpone    forces compilation of the next word, even if it
              is marked IMMEDIATE. Used inside definitions to
              defer execution of immediate words.

  value       (x --)
              Creates a named value that can be read by name
              and changed with TO.
              Usage: 42 value myval

  to          (x --)
              Change the value of a word created with VALUE.
              Usage: 99 to myval

  does>       ( -- )
              Defines the runtime behaviour of a CREATE'd word.
              Everything after DOES> runs when the created word
              is executed, with the data-field address on the
              stack.
              Usage: : array create cells allot does> swap
              cells + ;

  include     ( -- )
              Loads and runs a FORTH file.
              Usage: include maths.fs
\end{forthbox}

% ──────────────────────────────────────────────────────────────────
\subsection{Meta Programming}

\begin{forthbox}[Meta Programming]
  '        ( -- xt)
           Gets the execution token of the named word. Used to
           get a reference to a word for later execution.

  [']      ( -- xt)
           Compile-time version of ' (tick). Compiles the
           execution token of a word into the current definition.

  execute  (xt --)
           Executes the word whose execution token is on the
           stack. Enables dynamic word execution.

  literal  (x --)
           Compiles a numeric literal into the current word
           definition during compilation.

  state    ( -- addr)
           Returns the address of the STATE variable. STATE is 0
           when interpreting, -1 when compiling.

  [        (--)
           Switches from compilation mode to interpretation mode.

  ]        (--)
           Switches from interpretation mode to compilation mode.

  evaluate (c-addr u --)
           Interprets the string at c-addr of length u as if
           it were typed at the terminal.

  find     (c-addr -- c-addr 0 | xt 1 | xt -1)
           Searches the dictionary for the counted string at
           c-addr. Returns xt and 1 (immediate) or -1 (normal),
           or c-addr and 0 if not found.

  >body    (xt -- a-addr)
           Returns the data-field address of a CREATE'd word
           given its execution token.

  >in      ( -- a-addr)
           Returns the address of the variable holding the
           current input parse position.

  >number  (ud1 c-addr1 u1 -- ud2 c-addr2 u2)
           Converts characters at c-addr to a number using
           the current BASE. Stops at the first non-convertible
           character.
\end{forthbox}

% ──────────────────────────────────────────────────────────────────
\subsection{Arithmetic Operators}

Arithmetic is done using reverse Polish notation, so the numbers are required on the stack before the arithmetic operator can be used.

\begin{forthbox}[Arithmetic Operators]
  +        (n1 n2 -- sum)
           Takes the top two items from the stack and returns
           their sum to the stack.

  -        (n1 n2 -- difference)
           Takes the top two items from the stack and returns
           their difference to the stack.

  *        (n1 n2 -- product)
           Takes the top two items from the stack and returns
           their product to the stack.

  /        (n1 n2 -- quotient)
           Takes the top two items from the stack and returns
           their quotient to the stack.

  */       (n1 n2 n3 -- result)   result = (n1*n2/n3)
           Takes the top three items from the stack. The first
           two are multiplied together, then divided by the last
           item. The result is appended to the stack.

  mod      (n1 n2 -- modulus)
           Takes two items from the stack, divides the first by
           the second and the remainder is added to the stack.

  /mod     (n1 n2 -- modulus quotient)
           Takes the last two items from the stack and returns
           their modulus and their quotient to the stack.

  abs      (n -- |n|)
           Takes the last value from the stack and returns its
           absolute value.

  negate   (n -- -n)
           Takes the last item from the stack and returns its
           value negated.

  max      (n1 n2 -- max)
           Takes the last two values from the stack and returns
           the greatest value to the stack.

  min      (n1 n2 -- min)
           Takes the last two values from the stack and returns
           the smallest value to the stack.

  1+       (n -- n+1)
           Takes the top item from the stack, adds one to it and
           returns the result to the stack.

  1-       (n -- n-1)
           Takes the top item from the stack, minuses one from it
           and returns the result to the stack.

  2+       (n -- n+2)
           Takes the top item from the stack, adds two to it and
           returns the result to the stack.

  2-       (n -- n-2)
           Takes the top item from the stack, minuses two from it
           and returns the result to the stack.

  2*       (n -- n*2)
           Takes the top item from the stack, multiplies it by
           two and returns the result to the stack.

  2/       (n -- n/2)
           Takes the top item from the stack, halves it and
           returns the result to the stack.

  sumStack (n1 ... n -- sum)
           Adds all items on the stack and appends the sum.

  */mod    (n1 n2 n3 -- remainder quotient)
           Multiplies n1 by n2, then divides by n3. Returns
           both the remainder and the quotient.
\end{forthbox}

% ──────────────────────────────────────────────────────────────────
\subsection{Double-Cell Arithmetic}

Double-cell words operate on 64-bit values represented as two stack items (high cell on top). These are required by the Forth-2012 CORE standard for intermediate precision in calculations.

\begin{forthbox}[Double-Cell Arithmetic]
  s>d      (n -- d)
           Sign-extends a single-cell number to a double-cell
           number. Required before FM/MOD or SM/REM.

  m*       (n1 n2 -- d)
           Signed multiply of n1 by n2, producing a double-cell
           result. No overflow possible.

  um*      (u1 u2 -- ud)
           Unsigned multiply of u1 by u2, producing an unsigned
           double-cell result.

  fm/mod   (d n1 -- n2 n3)
           Floored division of double-cell d by n1. The quotient
           n3 rounds toward negative infinity. The remainder n2
           has the same sign as the divisor.

  sm/rem   (d n1 -- n2 n3)
           Symmetric division of double-cell d by n1. The
           quotient n3 rounds toward zero. The remainder n2 has
           the same sign as the dividend.

  um/mod   (ud u1 -- u2 u3)
           Unsigned division of unsigned double-cell ud by u1.
           Returns remainder u2 and quotient u3.
\end{forthbox}

% ──────────────────────────────────────────────────────────────────
\subsection{Numeric Output}

\begin{forthbox}[Numeric Output]
  <#       (--)
           Initialises pictured numeric output conversion.

  #        (ud1 -- ud2)
           Converts one digit from the double number and adds it
           to the pictured numeric output buffer.

  #s       (ud -- 0 0)
           Converts all remaining digits from the double number
           until it becomes zero.

  #>       (xd -- c-addr u)
           Completes pictured numeric conversion and returns the
           address and length of the formatted string.

  hold     (char --)
           Adds a character to the beginning of the pictured
           numeric output buffer.

  sign     (n --)
           If the number is negative, adds a minus sign to the
           pictured numeric output.

  .r       (n1 n2 --)
           Displays a number right-aligned in a field of n2
           characters width.

  u.       (u --)
           Displays an unsigned number followed by a space.

  u.r      (u n --)
           Displays an unsigned number right-aligned in a field
           of n characters width.
\end{forthbox}

% ──────────────────────────────────────────────────────────────────
\subsection{Comparison Operators}

\begin{notebox}
\texttt{f} is a flag or bool \texttt{[true\,|\,false]}. Per the Forth-2012 standard, \texttt{TRUE} is \texttt{-1} (all bits set) and \texttt{FALSE} is \texttt{0}. Any non-zero value is considered true by \texttt{IF}, \texttt{AND}, \texttt{OR}, and other conditional words.
\end{notebox}

\begin{forthbox}[Comparison Operators]
  =    (n1 n2 -- f)
       Takes the top two values from the stack, checks if they
       are equal, then returns the flag to the stack.

  <>   (n1 n2 -- f)
       Takes the top two values from the stack, checks if they
       are NOT equal, then returns the flag to the stack.

  <    (n1 n2 -- f)
       Takes the top two values from the stack, checks if n1 is
       less than n2, then returns the flag to the stack.

  >    (n1 n2 -- f)
       Takes the top two values from the stack, checks if n1 is
       greater than n2, then returns the flag to the stack.

  &&   (n1 n2 -- f)
       Logical AND. Returns TRUE (-1) if both values are
       non-zero, FALSE (0) otherwise.

  ||   (n1 n2 -- f)
       Logical OR. Returns TRUE (-1) if either value is
       non-zero, FALSE (0) otherwise.

  0=   (n1 -- f)
       Takes the top value from the stack, checks if it is equal
       to zero, then returns the flag to the stack.

  0<   (n1 -- f)
       Takes the top value from the stack, checks if it is less
       than zero, then returns the flag to the stack.

  0>   (n1 -- f)
       Takes the top value from the stack, checks if it is
       greater than zero, then returns the flag to the stack.

  u<   (u1 u2 -- f)
       Unsigned comparison. Returns true if u1 is less than u2
       when both are treated as unsigned values.
\end{forthbox}

% ──────────────────────────────────────────────────────────────────
\subsection{Bitwise Operators}

\begin{forthbox}[Bitwise Operators]
  and     (n1 n2 -- AND)
          Takes the top two values from the stack and returns the
          result of their logical AND to the stack.

  or      (n1 n2 -- OR)
          Takes the top two values from the stack and returns the
          result of their logical OR to the stack.

  xor     (n1 n2 -- XOR)
          Takes the top two values from the stack and returns the
          result of their logical XOR to the stack.

  invert  (n -- ~n)
          Bitwise inversion. Flips every bit of the top value
          on the stack.

  rshift  (n1 n2 -- n3)
          Takes the top two values from the stack and returns the
          result of n1 bitwise right-shifted by n2.

  lshift  (n1 n2 -- n3)
          Takes the top two values from the stack and returns the
          result of n1 bitwise left-shifted by n2.
\end{forthbox}

% ──────────────────────────────────────────────────────────────────
\subsection{Stack Operators}

\begin{forthbox}[Stack Operators]
  .s            ( -- )
                Displays the current stack to the terminal.

  .             (n -- )
                Takes the last value from the stack and pops it
                to terminal.

  dup           (n -- n n)
                The top value is taken from the stack, duplicated
                and both the original and copy returned.

  ?dup          (n -- 0 | n n)
                Duplicates the top item on the stack only if it
                is non-zero.

  nip           (n1 n2 -- n2)
                The first item below the top of the stack is
                dropped.

  swap          (n1 n2 -- n2 n1)
                The last two values are taken from the stack,
                swapped and then returned to the stack.

  drop          (n -- )
                Discards the top item from the stack.

  2dup          (n1 n2 -- n1 n2 n1 n2)
                The top two values from the stack are taken,
                copied and both original and copy returned.

  over          (n1 n2 -- n1 n2 n1)
                The second value from the stack is copied and
                added to the stack.

  pick          (nx ... n x -- nx ... n nx)
                Takes the top value from the stack and copies the
                nth value, appending it to the end of the stack.

  rot           (n1 n2 n3 -- n2 n3 n1)
                The top three values of the stack are taken and
                rotated left then returned to the stack.

  -rot          (n1 n2 n3 -- n3 n1 n2)
                The top three values of the stack are taken and
                rotated right then returned to the stack.

  tuck          (n1 n2 -- n2 n1 n2)
                The top two values from the stack are taken. The
                top value is copied and placed before the second
                value in the stack (third value deep).

  2drop         (n1 n2 -- )
                The top two values are taken from the stack and
                discarded.

  2swap         (n1 n2 n3 n4 -- n3 n4 n1 n2)
                The top four values are taken from the stack, and
                the order of the two pairs is reversed.

  2over         (n1 n2 n3 n4 -- n1 n2 n3 n4 n1 n2)
                The first four values are taken from the stack,
                with items 3 and 4 deep being copied and added.

  reverseStack  (n1 n2 ... nlast -- nlast ... n2 n1)
                The first shall become last and the last shall
                become first. The stack's order is reversed.

  depth         ( -- n)
                Returns the number of items currently on the data
                stack before depth was executed.

  roll          (xu xu-1 ... x0 u -- xu-1 ... x0 xu)
                Rotates u+1 items on the stack. The u-th item is
                moved to the top of the stack.

  stackCount    ( -- n)
                The number of items on the stack is added to the
                stack.
\end{forthbox}

% ──────────────────────────────────────────────────────────────────
\subsection{Return Stack Operators}

\begin{forthbox}[Return Stack Operators]
  >r   (n -- )
       Takes top value from the parameter stack and moves it to
       the return stack.

  r>   ( -- n)
       Takes top value from the return stack and moves it to the
       parameter stack.

  r@   ( -- n)
       Copies last item from return stack.

  i    ( -- n)
       Copies top item in return stack.

  i'   ( -- n)
       Copies second item in return stack.

  j    ( -- n)
       Copies third item in return stack.

  nth  ( -- n)
       Copies the nth item in the return stack.
\end{forthbox}

% ──────────────────────────────────────────────────────────────────
\subsection{Decision Operators}

\begin{notebox}
\texttt{f} is a flag or bool \texttt{[true\,|\,false]}.
\end{notebox}

\begin{forthbox}[Decision Operators]
  if     (f -- )
         Takes a flag from the stack and executes words if true
         flag.

  else   (f -- )
         If flag taken from if was false these words will be
         executed. This is optional.

  then   (f -- )
         Denotes the end of the if statement.

  not    (f -- f)
         Takes a flag value and reverses its result.

  within (n1 n2 n3 -- f)
         Takes three values from the stack and checks if n2 is a
         number within n1 and n3.

  true   ( -- -1)
         Pushes the true flag (-1, all bits set) to the stack.
         Per Forth-2012 standard.

  false  ( -- 0)
         Pushes the false flag (0) to the stack.

  case   ( -- )
         Begins a CASE...OF...ENDOF...ENDCASE switch statement.
         The selector value must be on the stack.

  of     (x1 x2 -- | x1)
         Compares the top of stack with the selector. If equal,
         drops both and executes the clause. If not, skips to
         ENDOF.

  endof  ( -- )
         Ends an OF clause. Branches to after ENDCASE.

  endcase (x -- )
         Ends a CASE statement. Drops the selector value if no
         OF clause matched.

  catch  (xt -- exception# | 0)
         Executes the word referred to by xt. If THROW is called
         during execution, catches the exception and returns
         the error code. Returns 0 on success.

  throw  (0 | exception# -- )
         Throws an exception. If 0, does nothing. A non-zero
         value unwinds the stack to the nearest CATCH.
\end{forthbox}

% ──────────────────────────────────────────────────────────────────
\subsection{Loop Operators}

\begin{forthbox}[Loop Operators]
  do     (n1 n2 --)
         n1 is endValue, n2 is startValue. Repeats words between
         do and loop from startValue until endValue is met.

  loop   ( -- )
         Ends loop.

  +loop  (n1 -- )
         Loops in n1 increments.

  every  (n1 -- )
         Performs word after every, every n1 seconds.

  stop   ( -- )
         Stops every running.

  begin  ( f -- )
         Marks the start of a while loop.

  while  ( f -- )
         Tests a condition within a BEGIN...REPEAT loop. If the
         condition is false, exits the loop.

  repeat ( -- )
         Marks the end of a BEGIN...WHILE...REPEAT loop and jumps
         back to BEGIN.

  until  ( -- )
         Ends BEGIN loop when condition is met.

  leave  ( -- )
         Exits the current DO...LOOP or DO...+LOOP immediately.
         Must be used within a loop.

  unloop ( -- )
         Discards the loop control parameters from the return
         stack. Used when exiting a loop early with LEAVE.

  ?do    (n1 n2 --)
         Like DO but skips the loop entirely if the start and
         end values are equal. Executes loop from n1 to n2-1.

  again  ( -- )
         Unconditional branch back to BEGIN. Creates an infinite
         loop (BEGIN...AGAIN). Use LEAVE or the Escape key to
         exit.
\end{forthbox}

% ──────────────────────────────────────────────────────────────────
\subsection{Base Operations}

Much like FORTH, Henceforth accepts integers as binary, octal, decimal and hexadecimal. Just be sure to change the base.

\begin{forthbox}[Base Operations]
  base      ( -- addr)
            Returns the address of the BASE variable (Forth-2012
            compliant). Use BASE @ to read and n BASE ! to write.
            Supports bases 2 through 36.

  decimal   (--)
            Sets BASE to 10.

  binary    (--)
            Sets BASE to 2.

  octal     (--)
            Sets BASE to 8.

  hex       (--)
            Sets BASE to 16.
\end{forthbox}

% ──────────────────────────────────────────────────────────────────
\subsection{Pop-up Operations}

\begin{forthbox}[Pop-up Operations]
  lineChart    (--)
               Displays a line chart representing the stack.

  barChart     (--)
               Displays a bar chart representing the stack.

  pieChart     (--)
               Displays a pie chart representing the stack with
               slices in proportion to its value's weight compared
               to the stack's sum.

  applemap     ( -- )
               Presents Apple Maps view.
\end{forthbox}

% ──────────────────────────────────────────────────────────────────
\subsection{Constants and Variables}

\begin{forthbox}[Constants and Variables]
  constant (n --)
           Takes a value from the stack and a name and creates a
           constant of that value called by the name provided.

  variable (--)
           Creates a variable whose name will be the first string
           of characters following variable, initially zero.

  !        (n address --)
           Stores a number to the address provided.

  @        (address -- n)
           Gives the value from address.

  ?        (address --)
           Prints the value from address to the terminal.

  c!       (char addr --)
           Stores a single character at the specified memory
           address.

  c@       (addr -- char)
           Fetches a single character from the specified memory
           address and places it on the stack.

  create   (--)
           Creates an array named by the first string of
           characters following the create word.

  ,        (n --)
           Takes the top two values from the stack and stores the
           second value as part of the array in the address of
           the first value.

  allot    (n --)
           Allocates n cells of data space. Used after CREATE to
           allocate memory for arrays or data structures.

  here     ( -- addr)
           Returns the address of the next available data space
           location.

  cells    (n1 -- n2)
           Converts a count to an address offset by multiplying
           by the cell size.

  cell+    (addr1 -- addr2)
           Adds the size of one cell to an address, advancing it
           to the next cell boundary.

  +!       (n addr --)
           Adds n to the value stored at the given address.

  count    (c-addr1 -- c-addr2 u)
           Converts a counted string (length byte followed by
           characters) to an address and length suitable for TYPE.

  bl       ( -- char)
           Pushes the ASCII value for space character (32) onto
           the stack.

  2@       (addr -- x1 x2)
           Fetches two consecutive cells from memory at the given
           address.

  2!       (x1 x2 addr --)
           Stores two consecutive values to memory at the given
           address.

  fill     (addr u char --)
           Fills u bytes of memory starting at addr with the
           specified character value.

  move     (addr1 addr2 u --)
           Copies u address units from addr1 to addr2. Handles
           overlapping memory regions correctly.

  erase    (addr u --)
           Fills u address units with zeros, effectively clearing
           the memory region.

  allocate (u -- a-addr ior)
           Allocates u cells of contiguous data space. Returns the
           start address and an I/O result code (0 = success, -1 =
           failure). Reuses previously freed blocks when possible.

  free     (a-addr -- ior)
           Releases previously allocated memory starting at addr.
           Returns 0 on success. Adjacent freed blocks are merged
           to reduce fragmentation.

  resize   (a-addr u -- a-addr2 ior)
           Changes the size of an allocation. If the new size fits
           within the existing block, returns the same address.
           Otherwise allocates a new block, copies existing data,
           and frees the old block.

  c,       (char --)
           Reserves one character of data space and stores char
           in it.

  char+    (c-addr1 -- c-addr2)
           Adds the size of one character to the address.

  chars    (n1 -- n2)
           Returns the size in address units of n characters.
           In this system, 1 char = 1 unit.

  align    ( -- )
           Aligns the data-space pointer to a cell boundary.

  aligned  (addr -- a-addr)
           Returns the first aligned address greater than or
           equal to addr.
\end{forthbox}

% ──────────────────────────────────────────────────────────────────
\subsection{Terminal Operators}

\begin{forthbox}[Terminal Operators]
  cr           (--)
               Newline added to terminal.

  space        (--)
               Adds a single space to terminal.

  spaces       (n --)
               Takes a value from the stack and adds that number
               of spaces to the terminal.

  emit         (n --)
               Takes a value from the stack and emits its ASCII
               code.

  (            (--)
               Starts inline comment. Everything between ( and )
               will be ignored. Can appear anywhere on a line.

  )            (--)
               Ends inline comment.

  \textbackslash            (--)
               Line comment. Everything after \textbackslash{} to
               the end of the line is ignored. Supported in .fs
               files. Example: \textbackslash{} this is a comment

  .s           (--)
               Show stack on terminal.

  .rs          (--)
               Shows return stack.

  clear        (--)
               Clears stack and return stack.

  bye          (--)
               Clears everything, resets to default settings.

  ."           (--)
               Start print. Everything between quotes is printed.

  "            (--)
               Ends print.

  s"           ( -- addr u)
               Creates a string literal. Returns the address and
               length of the string on the stack.

  type         (addr u --)
               Takes an address and length from the stack and
               displays the string at that address to the terminal.

  page         (--)
               Adds 24 new lines to the terminal.

  char         ( -- char)
               Parses the next space-delimited word and returns
               the ASCII value of its first character.

  [char]       ( -- n)
               Compile-time version of CHAR. Compiles the
               character value into the current definition.

  word         (char -- c-addr)
               Parses the input stream until the delimiter
               character is found and returns a counted string
               address.

  parse        (char -- c-addr u)
               Parses the input stream until the delimiter
               character is found and returns address and length.

  key          ( -- n)
               Waits for a single keypress and pushes its ASCII
               value onto the stack. Blocks execution until a key
               is pressed. Works with both hardware and software
               keyboards.

  ?terminal    ( -- flag)
               Returns true if the keyboard buffer has input or
               if a terminal interrupt has been requested (break
               key pressed), false otherwise. FIG-FORTH semantics.

  accept       (c-addr n1 -- n2)
               Receives up to n1 characters from the input source
               and stores them at the buffer address c-addr.
               Returns n2, the actual number of characters
               received.

  see          (--)
               Checks dictionary for word and outputs definition.

  sessionInput (--)
               Displays session input.

  dictionary   (--)
               Shows the user's defined words, constants and
               variables.

  vlist        (--)
               Displays a list of all the built-in words
               Henceforth offers.

  forth        (--)
               Resets vocabulary to standard built-in words,
               removing all user definitions.

  c"           ( -- c-addr)
               Creates a counted string literal. Returns the
               address of a counted string (length byte followed
               by characters).

  abort        (--)
               Clears all stacks and resets the interpreter state.

  abort"       (flag --)
               If the flag is non-zero, displays the message and
               ABORTs. Otherwise discards the message and
               continues.
               Usage: 0= abort" value must not be zero"

  quit         (--)
               Clears the return stack and enters interpretation
               state.

  source       ( -- c-addr u)
               Returns the address and length of the current
               input buffer.

  environment? (c-addr u -- false | i*x true)
               Queries the system for an environmental parameter.
               Returns false if unknown, or the value and true
               if known.
\end{forthbox}

% ──────────────────────────────────────────────────────────────────
\subsection{Other Words}

\begin{forthbox}[Other Words]
  date&time (-- second minute hour day month year)
            Adds the device's current time to the stack.

  pay       (--)
            Presents pay sheet, where you can send bitcoin.

  .fs       (--)
            Filename followed by .fs extension will load that
            file's contents to Henceforth.

  startup.fs
            A special file that runs automatically when
            Henceforth launches. Created as a blank template
            on first launch. Add your custom word definitions,
            aliases, and INCLUDE calls here so they are
            available every session.
\end{forthbox}

% ──────────────────────────────────────────────────────────────────
\subsection{Data Conversion}

Henceforth has two stacks: the integer stack (main) and the data stack (for raw bytes, used by Bitcoin Script operations). These words bridge between them.

\begin{forthbox}[Data Conversion]
  hex2data    ( -- )
              Converts the following hex string to bytes and
              pushes them onto the data stack.
              Usage: hex2data 48656c6c6f

  data2hex    ( -- )
              Converts the top data stack item to a hex string
              and displays it on the terminal.

  .ds         ( -- )
              Displays the contents of the data stack.

  >data       ( -- )
              Converts the following token to raw data bytes and
              pushes them onto the data stack.

  >addr       (n -- )
              Marks the top-of-stack integer value as a Bitcoin
              address string for transaction building.
\end{forthbox}

% ──────────────────────────────────────────────────────────────────
\subsection{String Operations}

\begin{forthbox}[String Operations]
  cmove       (c-addr1 c-addr2 u -- )
              Copies u characters from addr1 to addr2, proceeding
              low address to high (safe when addr2 > addr1).

  cmove>      (c-addr1 c-addr2 u -- )
              Copies u characters from addr1 to addr2, proceeding
              high address to low (safe for overlapping regions
              when addr2 < addr1).

  /string     (c-addr1 u1 n -- c-addr2 u2)
              Adjusts a string address and length by n characters.
              Advances the address by n and reduces the length by
              n.

  search      (c-addr1 u1 c-addr2 u2 -- c-addr3 u3 flag)
              Searches for string2 within string1. If found,
              returns the match address, remaining length, and
              true. If not found, returns the original string and
              false.

  blank       (c-addr u -- )
              Fills u characters of memory starting at addr with
              spaces (BL, ASCII 32).
\end{forthbox}

% ──────────────────────────────────────────────────────────────────
\subsection{Networking Operations}

\begin{forthbox}[Networking Operations]
  wocbitcoinprice         ( -- n)
                          Fetches the bitcoin price in USD and
                          adds it to the stack in pennies.

  woccirculatingsupply    (--)
                          Fetches the circulating supply of
                          bitcoin.

  wocblockchaininfo       (--)
                          Fetches latest general block info from
                          WhatsOnChain and prints to terminal.

  wocaddressinfo          (--)
                          Takes an address and displays its
                          information to the terminal.

  wocgetbalance           (--)
                          Takes an address and displays its
                          balance information to the terminal.

  wocgethistory           (--)
                          Takes an address and displays its
                          history information to the terminal.

  wocgetblockby           (n--)
                          Takes a value from the stack and
                          searches WhatsOnChain for info on that
                          block.

  wocgorillapoolfeequote  (--)
                          Fetches GorillaPool's fee quote via
                          ARC policy endpoint.

  taalfeequote            (--)
                          Fetches TAAL's fee quote via ARC
                          policy endpoint.
\end{forthbox}

% ──────────────────────────────────────────────────────────────────
\subsection{Transaction Builder}

The transaction builder words allow programmatic construction of Bitcoin transactions directly from the FORTH terminal. Build a transaction step by step, then broadcast it.

\begin{forthbox}[Transaction Builder]
  tx-new          ( -- )
                  Initialises a new transaction builder. Clears
                  any previous transaction state.

  tx-set-fee      (n -- )
                  Sets the fee rate in satoshis per kilobyte for
                  the current transaction.
                  Usage: 500 tx-set-fee

  tx-add-output   (n -- )
                  Adds a P2PKH output to the transaction. Expects
                  the amount in satoshis on the stack and the
                  destination address as the next word.
                  Usage: 1000 tx-add-output 1A1zP1...

  tx-add-data     ( -- )
                  Adds an OP\_RETURN data output to the current
                  transaction.

  tx-preview-on   ( -- )
                  Enables preview mode --- TX-BROADCAST will show
                  the transaction without actually broadcasting.

  tx-preview-off  ( -- )
                  Disables preview mode --- TX-BROADCAST will
                  broadcast to the network.

  tx-show         ( -- )
                  Displays the current transaction being built
                  (inputs, outputs, fee).

  tx-clear        ( -- )
                  Clears the transaction builder, discarding all
                  inputs and outputs.

  tx-broadcast    ( -- )
                  Builds, signs, and broadcasts the current
                  transaction to the BSV network via ARC.

  send            (n -- )
                  Standard payment in a single word. Reads the
                  destination address(es) from the last ." string
                  and the amount in satoshis from the stack.
                  Auto-selects UTXOs from the default wallet,
                  uses the wallet's fee rate, and broadcasts
                  via ARC.
                  Usage: ." 1A1zP1..." 1000 send
\end{forthbox}

\begin{notebox}
\textbf{Transaction Builder Example:} Send 1000 satoshis to an address with a custom fee rate using the step-by-step builder:
\begin{lstlisting}
tx-new
250 tx-set-fee
." 1ExampleAddress..." 1000 tx-add-output
tx-show
tx-broadcast
\end{lstlisting}
\end{notebox}

\begin{notebox}
\textbf{Standard Payment Example:} The \texttt{send} word wraps the entire PayView payment flow into a single command. It uses the \emph{default wallet} (marked with a star on the wallet card):
\begin{lstlisting}
\ Single recipient:
." 1A1zP1eP5QGefi2DMPTfTL5SLmv7DivfNa" 1000 send

\ Multiple recipients (space-separated):
." 1A1zP1... 1BvBMS..." 5000 send
\end{lstlisting}
Addresses are wrapped in \texttt{."} (dot-quote), which is FORTH's string literal syntax. The address string is placed on the terminal output, and \texttt{send} reads it from there. Multiple addresses can be space-separated within a single \texttt{."} string for multi-output transactions.
\end{notebox}

\subsubsection{Standard Transaction Script}

Every installation includes a \texttt{standard-transaction.fs} file in the Files tab. This is a Bitcoin Script file that defines the OP\_RETURN output attached to your payments. It is automatically loaded each time you open the Pay view.

\begin{lstlisting}
( Standard Transaction Script )

SCRIPT-BEGIN
OP_FALSE OP_RETURN
." Hello, you are welcome." >data data>script
SCRIPT-END
\end{lstlisting}

\textbf{How it works:}
\begin{enumerate}
\item \texttt{SCRIPT-BEGIN} activates script recording mode
\item \texttt{OP\_FALSE OP\_RETURN} emits the opcodes for a data-carrying output
\item \texttt{."} places the message on the terminal output
\item \texttt{>data} converts it to raw bytes on the data stack
\item \texttt{data>script} emits those bytes as \texttt{PUSHDATA} into the script buffer
\item \texttt{SCRIPT-END} finalises the compiled script
\end{enumerate}

To customise your standard transaction, edit the message string in the file. The Pay view re-reads the file from disk each time it opens, so edits take effect immediately. The file also appears in the Script tab of the script picker if you need to switch between scripts.

\begin{notebox}
\textbf{Note on comments:} Script files use \texttt{( )} parenthetical comments rather than \texttt{\textbackslash} line comments. The script picker flattens newlines when compiling, which causes \texttt{\textbackslash} comments to consume the rest of the file.
\end{notebox}

% ──────────────────────────────────────────────────────────────────
\subsection{Bitcoin Cryptography}

\begin{forthbox}[Bitcoin Cryptography]
  sign-msg     ( -- )
               Signs a message with the current wallet key using
               Bitcoin Signed Message (BSM).
               Usage: sign-msg <message>

  verify-msg   ( -- flag)
               Verifies a BSM signature against an address. Pushes
               true (-1) or false (0) to stack.
               Usage: verify-msg <address> <signature> <message>

  encrypt-msg  ( -- )
               Encrypts a message for a recipient using BRC-2
               Electrum ECIES (ECDH + AES-CBC + HMAC-SHA256).
               Wire-compatible with @bsv/sdk EncryptedMessage.
               Usage: encrypt-msg <recipientPubKeyHex> <message>

  decrypt-msg  ( -- )
               Decrypts a BRC-2 ECIES-encrypted message using the
               current wallet key. The plaintext is displayed.
               Usage: decrypt-msg <ciphertextHex>
\end{forthbox}

% ──────────────────────────────────────────────────────────────────
\subsection{BAP Identity}

BAP (Bitcoin Attestation Protocol) provides a decentralised identity system built on Bitcoin. Identity keys are derived using Type42 (BRC-42) key derivation from the wallet's master key.

\begin{forthbox}[BAP Identity Operations]
  bap-id       ( -- )
               Shows the BAP identity key (ID) and signing address
               for the current wallet.

  bap-sign     ( -- )
               Signs a message with the BAP identity key (Type42
               derived). Usage: bap-sign <message>

  bap-verify   ( -- flag)
               Verifies a BAP-signed message. Pushes true (-1) or
               false (0) to the stack.
               Usage: bap-verify <address> <signature> <message>

  bap-register ( -- )
               Builds a BAP ID registration transaction
               (OP_RETURN with BAP + AIP attestation).
\end{forthbox}

% ──────────────────────────────────────────────────────────────────
\subsection{Script Recording}

Script recording is a macro assembler for Bitcoin Script. When recording is active, \texttt{op\_*} words emit their opcodes into a script buffer rather than executing. Standard FORTH control flow still runs normally, allowing programmatic script construction.

\begin{forthbox}[Script Recording Operations]
  script-begin  ( -- )
                Start recording Bitcoin Script. op_* words emit
                opcodes to the script buffer. FORTH control flow
                still runs normally.

  script-end    ( -- )
                Stop recording and finalise the script. Displays
                the compiled script hex.

  script-show   ( -- )
                Display the current script buffer contents as hex.

  script-clear  ( -- )
                Discard the current script buffer and stop
                recording.

  script-use    ( -- )
                Set the recorded script as the selected transaction
                script (for building transactions).

  script-save   ( -- )
                Save the current script with a name.
                Usage: script-save <name>

  script-load   ( -- )
                Load a previously saved script by name.
                Usage: script-load <name>

  script-list   ( -- )
                List all saved scripts.
\end{forthbox}

% ──────────────────────────────────────────────────────────────────
\subsection{Script Bridge Words}

The bridge words connect FORTH's stacks to the script recording buffer, enabling dynamic script construction with computed values, strings, and file data.

\begin{forthbox}[Script Bridge Operations]
  >script       (n -- )
                Takes an integer from the stack and emits it as a
                Bitcoin Script number to the script buffer.

  data>script   ( -- )
                Takes the top item from the data stack and emits
                it as PUSHDATA to the script buffer. Use after
                crypto operations to inject computed hashes.

  text>script   ( -- )
                Converts the following token to UTF-8 bytes and
                emits as PUSHDATA to the script buffer.
                Usage: text>script image/jpeg

  file>script   ( -- )
                Reads a file's raw contents and emits them as
                PUSHDATA to the script buffer. For uploading
                images, videos, or data on-chain.
                Usage: file>script photo.jpg
\end{forthbox}

% ──────────────────────────────────────────────────────────────────
\subsection{Script Helper Words}

Pre-built script templates for common Bitcoin transaction patterns.

\begin{forthbox}[Script Helpers]
  make-p2pkh-script     ( -- )
                        Creates a standard P2PKH (Pay-to-Public-
                        Key-Hash) locking script from an address.

  make-opreturn-script  ( -- )
                        Creates an OP\_RETURN output script from
                        data on the stack.

  make-2of3-multisig    ( -- )
                        Creates a 2-of-3 multisig script from
                        three public keys on the stack.

  script-pushdata       ( -- )
                        Pushes raw data bytes to the script
                        recording buffer using OP\_PUSHDATA.
\end{forthbox}

\begin{notebox}
\textbf{Example:} Upload an image to the blockchain using the B:// protocol:
\begin{lstlisting}
script-begin
  op_false op_return
  text>script 19HxigV4QyBv3tHpQVcUEQyq1pzZVdoAut
  file>script photo.jpg
  text>script image/jpeg
  text>script binary
  text>script my_photo
script-end
script-use
pay
\end{lstlisting}
\end{notebox}

\begin{notebox}
\textbf{Example:} Build a hash-locked script with a computed HASH160:
\begin{lstlisting}
script-begin
  op_dup op_hash160
  " secret" op_sha256 op_hash160
  data>script
  op_equalverify op_checksig
script-end
\end{lstlisting}
Here \texttt{data>script} bridges the data stack (where the hash was computed) into the script buffer --- something impossible without the bridge words.
\end{notebox}

\begin{notebox}
\textbf{Example:} Use a FORTH loop to generate repetitive script patterns:
\begin{lstlisting}
script-begin
  10 0 do  op_1 op_add  loop
script-end
\end{lstlisting}
This emits 10 copies of \texttt{OP\_1 OP\_ADD} into the script. FORTH's full control flow (loops, variables, conditions) runs normally during recording --- only \texttt{op\_*} words emit to the buffer.
\end{notebox}

\begin{notebox}
\textbf{Example:} Build a simple P2PKH locking script:
\begin{lstlisting}
script-begin op_dup op_hash160 <pubKeyHash> op_equalverify op_checksig script-end
\end{lstlisting}
Scripts saved with \texttt{script-save} appear in the Files tab with an orange icon and can be selected as transaction scripts via \texttt{script-use} or the Script pop-up.
\end{notebox}

% ═════════════════════════════════════════════════════════════════
\newpage
\subsection{OPCodes}
% ═════════════════════════════════════════════════════════════════

\begin{forthbox}[Push Data OP Codes]
  op_0          ( -- 0)      Zero is pushed on to the stack.
  op_1          ( -- 1)      One is pushed on to the stack.
  op_2          ( -- 2)      Two is pushed on to the stack.
  op_3          ( -- 3)      Three is pushed on to the stack.
  op_4          ( -- 4)      Four is pushed on to the stack.
  op_5          ( -- 5)      Five is pushed on to the stack.
  op_6          ( -- 6)      Six is pushed on to the stack.
  op_7          ( -- 7)      Seven is pushed on to the stack.
  op_8          ( -- 8)      Eight is pushed on to the stack.
  op_9          ( -- 9)      Nine is pushed on to the stack.
  op_10         ( -- 10)     Ten is pushed on to the stack.
  op_11         ( -- 11)     Eleven is pushed on to the stack.
  op_12         ( -- 12)     Twelve is pushed on to the stack.
  op_13         ( -- 13)     Thirteen is pushed on to the stack.
  op_14         ( -- 14)     Fourteen is pushed on to the stack.
  op_15         ( -- 15)     Fifteen is pushed on to the stack.
  op_16         ( -- 16)     Sixteen is pushed on to the stack.

  op_false      ( -- 0)      Zero (false flag) pushed to stack.

  op_pushdata1  (--)   The byte after will contain how many
                       bytes are to be pushed on to the stack.

  op_pushdata2  (--)   The next two bytes will contain how many
                       bytes will be pushed on to the stack.

  op_pushdata4  (--)   The next four bytes will contain how many
                       bytes will be pushed on to the stack.

  op_1negate    ( -- -1)   Minus one pushed on to the stack.
\end{forthbox}

\begin{forthbox}[Control Flow OP Codes]
  op_nop       (--)
               Does nothing.

  op_ver       ( -- version)
               Pushes the transaction version onto the stack.
               Returns 1 (legacy) or 2 (Chronicle, post block
               943,816).

  op_if        (f--)
               Executes code depending on the expression's flag.

  op_notif     (n--)
               Executes code depending on the expression's flag.

  op_verif     ( threshold -- flag)
               Pushes TRUE if tx version >= threshold, FALSE
               otherwise. Chronicle upgrade.

  op_vernotif  ( threshold -- flag)
               Pushes TRUE if tx version < threshold, FALSE
               otherwise. Chronicle upgrade.

  op_else      (--)
               Executes if the preceding OP_IF, OP_NOTIF, or
               OP_ELSE were not executed.

  op_endif     (--)
               Ends the if/else statement.

  op_verify    (f--)
               Marks the transaction as INVALID if the flag is
               false.

  op_return    (--)
               Ends script. Often used to add data to the
               blockchain.
\end{forthbox}

\begin{forthbox}[Stack Operator OP Codes]
  op_toaltstack    (n--)
                   Takes the top value from the stack and pushes
                   it to the alt stack.

  op_fromaltstack  (--n)
                   Takes the top value from the alt stack and
                   pushes it to the stack.

  op_2drop         (n n1--)
                   Discards the top 2 elements from the stack.

  op_2dup          (n n1 -- n n1 n n1)
                   Duplicates the top two values from the stack
                   adding back both original and copy.

  op_3dup          (n n1 n2 -- n n1 n2 n n1 n2)
                   Duplicates the top three values from the stack
                   adding back both original and copy.

  op_2over         (n n1 n2 n3 -- n n1 n2 n3 n n1)
                   Copies the fourth and third items in the stack
                   over the first and second.

  op_2rot          (n n1 n2 n3 n4 n5 -- n2 n3 n4 n5 n n1)
                   Copies the fifth and sixth items in the stack
                   over the first and second.

  op_2swap         (n n1 n2 n3 -- n2 n3 n n1)
                   Takes four values from the stack and swaps
                   their order as pairs.

  op_ifdup         (n -- 0 | n n)
                   Duplicates the top value from the stack if it
                   is NOT zero. Leaves 0 unchanged.

  op_depth         (--n)
                   Adds the number of items on the stack to the
                   stack.

  op_drop          (n--)
                   Discards the top item from the stack.

  op_dup           (n -- n n)
                   Duplicates the top item from the stack.

  op_nip           (n n1 -- n1)
                   Discards the second item in the stack.

  op_over          (n n1 -- n n1 n)
                   Copies the second item in the stack over the
                   top.

  op_pick          (xu ... x0 u -- xu ... x0 xu)
                   Copies the u-th item on the stack to the top.

  op_roll          (xu xu-1 ... x0 u -- xu-1 ... x0 xu)
                   Moves the u-th item on the stack to the top,
                   shifting items above it down.

  op_rot           (n n1 n2 -- n1 n2 n)
                   The top three items on the stack are rotated
                   left, with the third item going to the top.

  op_swap          (n n1 -- n1 n)
                   The top two items on the stack's order is
                   reversed.

  op_tuck          (n n1 -- n1 n n1)
                   The top item on the stack is copied in front
                   of the second item.
\end{forthbox}

\begin{forthbox}[Size \& Separator OP Codes]
  op_size           (x -- x n)
                    Pushes the byte length of the top stack element
                    without consuming it.

  op_codeseparator  (--)
                    Marks the beginning of signature-checked data.
                    All signature checks only use data after the
                    most recent OP\_CODESEPARATOR.

  op_true           ( -- 1)
                    Pushes 1 onto the stack. Alias for OP\_1.
\end{forthbox}

\begin{forthbox}[Splice Operator OP Codes]
  op_cat      (--)
              Concatenates the two strings provided.

  op_split    (n n1--)
              Splits byte sequence n at n1.

  op_num2bin  (n n1--)
              Converts value n into a byte sequence of length n1.

  op_bin2num  (n--)
              Converts byte sequence n into a numerical value.

  op_substr   (start length --) data: (data -- substring)
              Extracts a substring by start index and length.
              Chronicle upgrade (0xb3).

  op_left     (n --) data: (data -- left)
              Extracts the leftmost N bytes.
              Chronicle upgrade (0xb4).

  op_right    (n --) data: (data -- right)
              Extracts the rightmost N bytes.
              Chronicle upgrade (0xb5).
\end{forthbox}

\begin{forthbox}[Bitwise Operator OP Codes]
  op_invert       (n -- n)
                  Flips all the bits in the input.

  op_and          (n1 n2 -- AND)
                  Logical AND performed and the result added to
                  the stack.

  op_or           (n1 n2 -- OR)
                  Logical OR performed and the result added to
                  the stack.

  op_xor          (n1 n2 -- XOR)
                  Logical XOR performed and the result added to
                  the stack.

  op_equal        (n n1 -- f)
                  Returns a flag depending on if n1 and n2 are
                  equal. 1 if true, 0 if false.

  op_equalverify  (n1 n2 -- f)
                  OP_EQUAL then OP_VERIFY.

  op_reserved1    (--)
                  Renders transaction invalid unless in
                  unexecuted OP_IF statement.

  op_reserved2    (--)
                  Renders transaction invalid unless in
                  unexecuted OP_IF statement.
\end{forthbox}

\begin{forthbox}[Arithmetic Operator OP Codes]
  op_1add                (n -- n+1)
                         1 is added to the top value on the stack.

  op_1sub                (n -- n-1)
                         1 is subtracted from the top value.

  op_2mul                (n -- n*2)
                         The top value on the stack is doubled.

  op_2div                (n -- n/2)
                         The top value on the stack is halved.

  op_negate              (n -- -n)
                         The top value's sign is flipped.

  op_abs                 (n -- |n|)
                         The absolute value of the top item on
                         the stack is returned.

  op_not                 (n -- f)
                         Logical NOT. Returns 1 if the input is 0,
                         otherwise returns 0.

  op_0notequal           (n -- f)
                         Returns true if the value is not zero,
                         else false.

  op_add                 (n n1 -- n3)
                         Top two items from the stack are summed
                         and the value returned.

  op_sub                 (n n1 -- n3)
                         The top value from the stack is
                         subtracted from the second value.

  op_mul                 (n n1 -- n3)
                         The top value is multiplied by the
                         second value in the stack.

  op_div                 (n n1 -- n3)
                         The second value in the stack is divided
                         by the top value.

  op_mod                 (n n1 -- n3)
                         The second value is divided by the top
                         value and the remainder returned.

  op_lshift              (n n1 -- n3)
                         The second value is logical left shifted
                         by the top value.

  op_rshift              (n n1 -- n3)
                         The second value is logical right shifted
                         by the top value.

  op_lshiftnum           (n shift -- result)
                         Numeric left shift preserving sign. Distinct
                         from op\_lshift which is bitwise on bytes.
                         Chronicle upgrade (0xb6).

  op_rshiftnum           (n shift -- result)
                         Numeric right shift preserving sign. Distinct
                         from op\_rshift which is bitwise on bytes.
                         Chronicle upgrade (0xb7).

  op_booland             (n n1 -- f)
                         True flag returned if both n and n1 are
                         non-zero, else false.

  op_boolor              (n n1 -- f)
                         True flag returned if either n or n1 is
                         not 0, else false.

  op_numequal            (n n1 -- f)
                         True if n is equivalent to n1, else
                         false.

  op_numequalverify      (n n1 -- f)
                         OP_NUMEQUAL then OP_VERIFY.

  op_numnotequal         (n n1 -- f)
                         True if n is NOT equal to n1, else
                         false.

  op_lessthan            (n n1 -- f)
                         True if n is less than n1, else false.

  op_greaterthan         (n n1 -- f)
                         True if n is greater than n1, else
                         false.

  op_lessthanorequal     (n n1 -- f)
                         True if n is less than or equal to n1,
                         else false.

  op_greaterthanorequal  (n n1 -- f)
                         True if n is greater than or equal to
                         n1, else false.

  op_min                 (n n1 -- n)
                         The top two values from the stack are
                         taken and the smaller returned.

  op_max                 (n n1 -- n)
                         The top two values from the stack are
                         taken and the larger returned.

  op_within              (n n1 n2 -- f)
                         True is returned if n is in between n1
                         and n2.
\end{forthbox}

\begin{forthbox}[Crypto Operator OP Codes]
  op_ripemd160             (--)
  op_sha1                  (--)
  op_sha256                (--)
  op_hash160               (--)
  op_hash256               (--)
  op_checksig              (--)
  op_checksigverify        (--)
  op_checkmultisig         (--)
  op_checkmultisigverify   (--)
\end{forthbox}

\begin{forthbox}[Lock Time Operator OP Codes]
  op_nop2           (--)
  op_nop3           (--)
  op_pubkeyhash     (--)
  op_pubkey         (--)
  op_invalidopcode  (--)
\end{forthbox}

\begin{forthbox}[Reserved OP Codes]
  op_nop1    (--)
  op_nop9    (--)
  op_nop10   (--)
\end{forthbox}

\begin{notebox}
\textbf{Chronicle Upgrade (Block 943,816):} NOP4--8 (0xb3--0xb7) were repurposed as \texttt{op\_substr}, \texttt{op\_left}, \texttt{op\_right}, \texttt{op\_lshiftnum}, and \texttt{op\_rshiftnum}. The transaction version (\texttt{op\_ver}) auto-switches from 1 to 2 at the activation block.
\end{notebox}

% ═════════════════════════════════════════════════════════════════
\chapter{Bitcoin Wallet}
% ═════════════════════════════════════════════════════════════════

Henceforth includes a full Bitcoin SV (BSV) wallet, enabling users to create wallets, receive and send transactions, and interact with the blockchain --- all from the same app that runs FORTH programs.

\section{Wallet Architecture}

The wallet is built on a hierarchical deterministic (HD) key structure, meaning a single seed phrase generates an unlimited number of addresses.

\begin{itemize}
\item \textbf{Seed Phrase:} A 12-word BIP-39 mnemonic is generated when creating a new wallet. This seed phrase is the master backup --- from it, every key can be regenerated.
\item \textbf{Key Derivation:} Two derivation methods are supported:
  \begin{itemize}
  \item \textbf{Type42 (BRC-42)} --- the default for new wallets. Uses elliptic curve key tweaking (private key addition, public key multiplication) for more flexible derivation paths.
  \item \textbf{BIP-44} --- the legacy standard. Retained for restoring wallets created by other software. Uses the standard path \texttt{m/44'/236'/0'/0/n}.
  \end{itemize}
\item \textbf{Address Rotation:} After each transaction, a new receiving address is derived automatically. This improves privacy by avoiding address reuse.
\end{itemize}

\begin{notebox}
\textbf{Security:} Private keys and seed phrases are stored exclusively in the iOS Keychain, protected by the device's Secure Enclave. They never leave the device and are never transmitted over the network.
\end{notebox}

\subsection{Address Discovery}

The wallet discovers addresses with funds using a shared Type42 scanner (\texttt{scanType42Addresses}). The scanner is \emph{derivationType-independent} --- it always tries Type42 first if the master key exists in Keychain, regardless of what \texttt{derivationType} says on the wallet card. This prevents funds from becoming invisible if the derivation type metadata is corrupted.

\begin{itemize}
\item \textbf{Dynamic scan range:} Scans at least 20 addresses (gap limit) past the last address with activity, with a minimum floor of 50 addresses.
\item \textbf{Receive and change chains:} Scans both \texttt{1-wallet-\{N\}} (receive) and \texttt{0-change-\{N\}} (change) invoice numbers.
\item \textbf{BIP-44 fallback:} After Type42, scans standard BIP-44 paths. A wallet can have funds on both derivation paths.
\item \textbf{Auto-repair:} If Type42 addresses are discovered but \texttt{derivationType} is wrong, the scanner automatically corrects it and persists the fix.
\end{itemize}

\section{Wallet Cards}

Each wallet is represented by a \emph{Wallet Card} in the app interface. A card stores:

\begin{itemize}
\item \textbf{Wallet title} --- a user-chosen name
\item \textbf{Main address} --- the primary receiving address
\item \textbf{Derived addresses} --- additional addresses generated through key rotation
\item \textbf{Balance} --- confirmed and unconfirmed satoshis, reconciled from the UTXO set
\item \textbf{Derivation type} --- Type42 or BIP-44
\end{itemize}

Multiple wallets can be created and managed simultaneously. The wallet card interface supports creating, renaming, and deleting wallets, with biometric (Face ID / Touch ID) authentication required to access sensitive details.

\subsubsection*{Default Wallet}

One wallet can be marked as the \emph{default wallet}, indicated by a star on the wallet card. The default wallet is used by the terminal \texttt{send} word and other non-UI operations. If no default is set, the first wallet is used.

To change the default wallet, open the wallet details view and tap \textbf{Set as Default Wallet}. The setting persists across app launches.

\section{Transactions}

\subsection{Receiving}

To receive BSV, share your wallet's receiving address or scan its QR code. The wallet monitors for incoming transactions using:

\begin{enumerate}
\item \textbf{JungleBus REST polling} --- checks every 30 seconds for new transactions via the \texttt{GET /v1/address/get/\{address\}} endpoint
\item \textbf{JungleBus WebSocket} --- optional real-time push notifications for near-instant detection (requires a subscription from the JungleBus dashboard)
\item \textbf{Periodic sync} --- full balance and UTXO refresh every 60 seconds as a safety net
\end{enumerate}

When a new transaction is detected, the wallet automatically refreshes the balance, UTXOs, and transaction history.

\subsection{Sending}

Transactions are built locally on the device:

\begin{enumerate}
\item \textbf{UTXO Selection:} The wallet selects unspent transaction outputs to fund the payment, avoiding ordinal-protected UTXOs.
\item \textbf{Fee Calculation:} Fees are calculated using the current rate from the preferred miner (GorillaPool or TAAL) via ARC policy endpoints.
\item \textbf{Change Output:} Any excess is returned to a fresh change address derived from the wallet.
\item \textbf{Signing:} Each input is signed with the correct private key for its address (Type42 or BIP-44 derived).
\item \textbf{Anti-Fee-Sniping:} The \texttt{nLockTime} field is set to the current block height, preventing miners from re-mining old transactions for higher fees.
\item \textbf{Broadcasting:} The signed transaction is submitted via ARC with automatic failover between GorillaPool and TAAL.
\end{enumerate}

The \texttt{pay} word in the FORTH terminal opens the transaction builder. The pay view has three tabs for selecting what to broadcast:

\begin{itemize}
\item \textbf{OP\_RETURN} --- write a text message to embed on-chain
\item \textbf{Script} --- select a saved script or a \texttt{.fs} script file
\item \textbf{Upload} --- pick a photo, video, or file from the device and upload it on-chain via the B:// protocol
\end{itemize}

\begin{notebox}
\textbf{Smart Sync:} When a payment is initiated, UTXOs are loaded from the local disk cache first (instant). A network fetch only occurs on genuine first launch when no UTXO file exists. Fee quotes are cached for 5 minutes. A times-table circle animation (cardioid pattern) plays during broadcast. The payment flow trusts local UTXO state --- if a UTXO turns out to be already spent, ARC will reject with \texttt{DOUBLE\_SPEND\_ATTEMPTED}.
\end{notebox}

After a successful broadcast, the wallet builds a local transaction model from the data it already has (inputs, outputs, addresses, amounts) and inserts it into the transaction history immediately --- no network fetch required. The ARC-returned \texttt{txid} is verified against the locally computed \texttt{txid} as cryptographic proof that the network accepted exactly what was built. Spent UTXOs are removed from memory (preventing double-spend on rapid re-send) and change UTXOs are added. Full transaction details from WhatsOnChain overwrite the local entry on the next sync.

\begin{notebox}
\textbf{Trust Local State:} Per the Bitcoin white paper (\S2), we built the transaction and signed it with our private keys --- we know exactly what it contains. There is no need to ask a third party (WhatsOnChain) what we just sent. ARC confirms the network saw it; the rest we already know.
\end{notebox}

\subsection{Transaction History}

The wallet displays a chronological transaction history with:

\begin{itemize}
\item \textbf{Direction detection} --- incoming (received) vs outgoing (sent), determined by checking whether any inputs belong to the wallet
\item \textbf{Grouping} --- transactions are grouped by: Pending, Today, Yesterday, This Week, This Month, Earlier
\item \textbf{Search and filter} --- filter by amount, date range, or search by transaction ID
\item \textbf{SPV verification} --- transactions can be verified against block headers using Simplified Payment Verification
\end{itemize}

\section{UTXOs and Balance}

A UTXO (Unspent Transaction Output) represents a discrete amount of Bitcoin that can be spent. The wallet tracks all UTXOs across all addresses.

\begin{notebox}
\textbf{Balance from UTXOs:} The displayed balance is computed solely from the UTXO set --- no separate balance API call is needed. When UTXOs are fetched, the balance is calculated as:
\begin{lstlisting}
confirmed   = sum of UTXOs with block height > 0
unconfirmed = sum of UTXOs with block height = 0
total       = confirmed + unconfirmed
\end{lstlisting}
Satoshi conversions use the named constant \texttt{satoshisPerBSV} (100,000,000) and price conversions use \texttt{priceCentsDivisor} (100.0). Raw magic numbers are never used.
\end{notebox}

\subsection{Ordinal Protection}

1Sat Ordinals are NFTs inscribed on single-satoshi UTXOs. The wallet protects ordinals at two levels:

\begin{enumerate}
\item \textbf{Transaction builder:} UTXOs with a value of 1 satoshi are excluded from UTXO selection during transaction building --- a fast local check with no network call.
\item \textbf{Inscription scanner:} For display in the Ordinals view, 1-sat UTXOs are scanned by fetching the raw transaction and detecting the \texttt{OP\_FALSE OP\_IF ... OP\_ENDIF} inscription pattern. Non-1-sat UTXOs are skipped entirely (ordinals are always exactly 1 sat). Results are cached per outpoint --- once a UTXO has been scanned, it is never re-fetched from the network.
\end{enumerate}

\section{API Providers}

The wallet uses a hybrid architecture with multiple blockchain service providers:

\begin{center}
\begin{tabular}{@{} l c c c c @{}}
\toprule
\textbf{Capability}      & \textbf{WhatsOnChain} & \textbf{GorillaPool} & \textbf{TAAL} & \textbf{JungleBus} \\
\midrule
UTXOs (balance derived)   & \cmark & \xmark & \xmark & \xmark \\
Transaction history       & \cmark & \xmark & \xmark & \cmark \\
SPV Merkle proofs         & \cmark & \xmark & \xmark & \xmark \\
BSV price                 & \cmark & \xmark & \xmark & \xmark \\
Fee quotes                & \xmark & \cmark & \cmark & \xmark \\
Broadcast (ARC)           & \xmark & \cmark & \cmark* & \xmark \\
Real-time monitoring      & \xmark & \xmark & \xmark & \cmark \\
\bottomrule
\end{tabular}

\noindent *TAAL requires a user-configured API key. GorillaPool is the primary broadcaster; TAAL is failover.
\end{center}

Each provider handles what it does best. API keys and subscription IDs are stored securely in the iOS Keychain and can be configured in Settings $\rightarrow$ API \& Miners.

\section{Security}

\begin{itemize}
\item \textbf{Biometric authentication} --- Face ID or Touch ID required to view private keys, seed phrases, and wallet details. Sessions expire after 120 seconds of inactivity.
\item \textbf{Keychain storage} --- All sensitive material (keys, seeds, API tokens) stored with \texttt{.whenUnlockedThisDeviceOnly} accessibility --- encrypted at rest, tied to the device.
\item \textbf{Seed verification} --- After creating a wallet, users are prompted to verify their seed phrase via a 3-word quiz.
\item \textbf{No server dependency} --- The wallet is fully self-custodial. No account, no registration, no server holds your keys.
\end{itemize}

\end{document}
